\documentclass[11pt]{article}

\usepackage[margin=1in]{geometry}
\usepackage[T1]{fontenc}
\usepackage[utf8]{inputenc}
\usepackage{lmodern}
\usepackage{microtype}
\usepackage{amsmath,amssymb}
\usepackage{hyperref}
\usepackage{enumitem}
\usepackage{xcolor}

\hypersetup{
  colorlinks=true,
  linkcolor=blue!50!black,
  citecolor=blue!50!black,
  urlcolor=blue!50!black
}

\title{Lean-Backed NS Interface Note}
\author{}
\date{\today}

\begin{document}
\maketitle

\begin{abstract}
This note records only the Navier--Stokes material currently backed by Lean in
the MeTTapedia repository.  It intentionally omits manuscript-specific claims
about Ben Goertzel's December 11, 2025 SG-Cole-Hopf draft unless they have been
encoded and checked in Lean.

The present Lean status is interface-level rather than proof-complete.  The
formalized material isolates reusable shells for cutoff continuity, topological
Cole--Hopf transport, approximation interfaces, finite identity-energy
algebra, abstract vorticity bounds, a kernel-semigroup wrapper, and two small
concrete models showing both positive approximation instantiations and wrong
topology obstructions, including one exact same-state-space split where a
summable geometric observable is continuous and a stronger tail detector is
not.  It also now contains an abstract good observable class, an abstract
tail-spike no-go theorem, and a first concrete soft-vs-hard chart fork on that
same shared mode space, together with more manuscript-shaped Sobolev-style and
energy-style soft-vs-hard forks on that same space, plus a log-energy fork
showing that a smoother radial profile still lands on the same good/bad split,
and finally a generic energy-profile theorem packaging that split for any
continuous radial profile dominated by the identity on nonnegative inputs.  It
does \emph{not} yet formalize the manuscript's specific chart, local
hypotheses \((H1)\), \((H2)\), or the BKM-closing analytic bridge.
\end{abstract}

\section{Scope}

The Navier--Stokes Lean work currently lives in:
\begin{enumerate}[leftmargin=1.5em]
  \item \texttt{FluidDynamics/NavierStokes/CutoffContinuity.lean},
  \item \texttt{FluidDynamics/NavierStokes/WeakTopologyCountermodel.lean},
  \item \texttt{FluidDynamics/NavierStokes/ProductTopologyL1Countermodel.lean},
  \item \texttt{FluidDynamics/NavierStokes/ColeHopfTopology.lean},
  \item \texttt{FluidDynamics/NavierStokes/ApproximationInterface.lean},
  \item \texttt{FluidDynamics/NavierStokes/ChartApproximationModel.lean},
  \item \texttt{FluidDynamics/NavierStokes/GeometricModeApproximation.lean},
  \item \texttt{FluidDynamics/NavierStokes/ModeStateProductTopologyObstruction.lean},
  \item \texttt{FluidDynamics/NavierStokes/ModeStateProductTopologyContinuity.lean},
  \item \texttt{FluidDynamics/NavierStokes/ModeStateObservableClass.lean},
  \item \texttt{FluidDynamics/NavierStokes/ModeStateTailSensitivityObstruction.lean},
  \item \texttt{FluidDynamics/NavierStokes/ModeStateChartFork.lean},
  \item \texttt{FluidDynamics/NavierStokes/ModeStateSobolevChartFork.lean},
  \item \texttt{FluidDynamics/NavierStokes/ModeStateEnergyChartFork.lean},
  \item \texttt{FluidDynamics/NavierStokes/ModeStateLogEnergyChartFork.lean},
  \item \texttt{FluidDynamics/NavierStokes/ModeStateEnergyProfileFork.lean},
  \item \texttt{FluidDynamics/NavierStokes/IdentityEnergy.lean},
  \item \texttt{FluidDynamics/NavierStokes/ColeHopfInterface.lean},
  \item \texttt{FluidDynamics/NavierStokes/KernelSemigroupInterface.lean}.
\end{enumerate}

These modules compile without \texttt{sorry}.  They provide a formal shell for
what a future SG-Cole-Hopf proof attempt would need to instantiate.

\section{Lean-Backed Results}

\subsection{Cutoff Continuity Shell}

In \texttt{CutoffContinuity.lean}, the generic cutoff-shaped potential
\[
\mathrm{cutoffPotential}(x) = \mathrm{cutoff}(\mathrm{radius}(x))
\cdot \mathrm{statistic}(x)
\]
is formalized together with:
\begin{enumerate}[leftmargin=1.5em]
  \item continuity of the cutoff potential when the scalar pieces are
  continuous;
  \item a sequential obstruction theorem: if the cutoff factor and statistic
  converge to the wrong product, then the cutoff potential cannot converge to
  its claimed value;
  \item the corresponding pointwise non-continuity criterion.
\end{enumerate}

So Lean already captures the generic shape of a chart-cutoff continuity
argument and its failure mode.

\subsection{Concrete Weak-Topology Countermodel}

In \texttt{WeakTopologyCountermodel.lean}, the abstract cutoff/Cole--Hopf shell
is instantiated on the concrete state space \(\mathbb{R}\times\mathbb{R}\)
equipped with the topology induced by the first coordinate.  Lean proves:
\begin{enumerate}[leftmargin=1.5em]
  \item a sequence can converge in the ambient topology while changing a hidden
  coordinate;
  \item a cutoff potential depending on that hidden coordinate is not
  continuous;
  \item the associated Cole--Hopf transform is likewise not continuous.
\end{enumerate}

This is a formal countermodel for the generic wrong-topology failure mode:
ambient convergence can be too weak for the cutoff datum it is supposed to
control.

\subsection{Infinite-Mode Product-Topology Countermodel}

In \texttt{ProductTopologyL1Countermodel.lean}, Lean moves the same mismatch to
the raw infinite-mode sequence space \(\mathbb{N}\to\mathbb{R}\) with the
product topology.  It proves:
\begin{enumerate}[leftmargin=1.5em]
  \item finite-mode truncations converge coordinatewise in the product
  topology;
  \item a tail spike moving to higher and higher coordinates converges to the
  zero sequence in that topology;
  \item the stronger tail-sensitive observable
  \(\sum_n |x_n|\) stays equal to \(1\) on that spike sequence and is therefore
  not continuous at zero;
  \item the associated Cole--Hopf datum is likewise not continuous.
\end{enumerate}

This is the first Lean-backed infinite-mode obstruction in the NS lane that
directly separates weak coordinatewise convergence from a stronger tail norm.

\subsection{Topological Cole--Hopf Shell}

In \texttt{ColeHopfTopology.lean}, Lean formalizes:
\begin{enumerate}[leftmargin=1.5em]
  \item continuity of the scalar map \(y \mapsto \exp(-y/(2\nu))\);
  \item continuity of the induced field-level Cole--Hopf transform;
  \item injectivity when \(\nu \neq 0\);
  \item transfer of potential-limit mismatch into Cole--Hopf-limit mismatch;
  \item the same obstruction specialized to cutoff-shaped potentials.
\end{enumerate}

This means the topological \(W \mapsto \phi\) step is already packaged in Lean,
provided one supplies continuity or convergence of the underlying potential.

\subsection{Approximation Interface}

In \texttt{ApproximationInterface.lean}, Lean packages an abstract positive
repair obligation:
\begin{quote}
if an approximation scheme makes the chart radius observable and the chart
statistic converge, then the cutoff potential and the Cole--Hopf datum also
converge.
\end{quote}

This is useful because it isolates exactly what an approximation/truncation
theorem would have to prove in order for the cutoff/Cole--Hopf layer to behave
well.

\subsection{Concrete Chart/Truncation Model}

In \texttt{ChartApproximationModel.lean}, the abstract approximation interface
is instantiated on the concrete chart state space \(\mathbb{R}\times\mathbb{R}\)
with:
\begin{enumerate}[leftmargin=1.5em]
  \item a hidden-coordinate radius observable;
  \item a simple mixed chart statistic;
  \item a truncation map that drops the hidden coordinate only at stage \(0\)
  and is exact thereafter.
\end{enumerate}

Lean then proves:
\begin{enumerate}[leftmargin=1.5em]
  \item convergence of the radius observable under truncation;
  \item convergence of the chart statistic under truncation;
  \item the resulting convergence of the cutoff potential;
  \item the resulting convergence of the induced Cole--Hopf datum.
\end{enumerate}

So the approximation interface is now backed by a concrete positive model, not
just an abstract shell.

\subsection{Geometric Finite-Mode Approximation Model}

In \texttt{GeometricModeApproximation.lean}, the same approximation interface is
instantiated on bounded infinite coefficient sequences.  Lean formalizes:
\begin{enumerate}[leftmargin=1.5em]
  \item a geometrically weighted mode radius;
  \item a chart statistic built from the zeroth visible mode plus that radius;
  \item finite-mode truncation keeping only the first \(N\) coefficients;
  \item summability of the weighted mode series under the uniform coefficient
  bound;
  \item convergence of the weighted radius under finite-mode truncation;
  \item convergence of the derived chart statistic, the cutoff potential, and
  the induced Cole--Hopf datum.
\end{enumerate}

This is the first genuinely infinite-mode positive instantiation in the NS Lean
lane: the approximation is not just ``wrong at stage \(0\), exact thereafter,''
but a real truncation-to-limit model on an infinite sequence state space.

\subsection{Same-State-Space Product-Topology Obstruction}

In \texttt{ModeStateProductTopologyObstruction.lean}, Lean returns to the exact
\texttt{ModeState} and \texttt{truncateModes} objects from
\texttt{GeometricModeApproximation.lean}, but equips that same state space with
the coordinatewise product topology induced by the coefficient map.  It proves:
\begin{enumerate}[leftmargin=1.5em]
  \item the existing finite-mode truncations \(\texttt{truncateModes}\) do
  converge to the target state in this product topology;
  \item a spike sequence still converges to the zero state in that same
  topology;
  \item a stronger tail-presence observable, which records whether any
  nonzero coefficient exists, is not continuous there;
  \item the associated Cole--Hopf datum is likewise not continuous.
\end{enumerate}

So Lean now contains a positive and a negative result on one common
infinite-mode object: the geometric weighted radius behaves well under
truncation, while a stronger tail-sensitive observable on the same state space
does not.

\subsection{Same-State-Space Product-Topology Continuity}

In \texttt{ModeStateProductTopologyContinuity.lean}, Lean proves the matching
positive theorem on the exact same \texttt{ModeState} equipped with the exact
same coordinatewise product topology:
\begin{enumerate}[leftmargin=1.5em]
  \item each coordinate observable \(x \mapsto x_n\) is continuous;
  \item each weighted mode term \(x \mapsto |x_n| 2^{-n}\) is continuous;
  \item the summable geometric radius \(\sum_n |x_n| 2^{-n}\) is continuous;
  \item the derived chart statistic, cutoff potential, and Cole--Hopf datum are
  therefore continuous as well.
\end{enumerate}

So the current NS Lean lane now cleanly separates two observables on one common
infinite-mode state space: a summably dominated geometric radius that behaves
well in the product topology, and a stronger tail detector that does not.

\subsection{Abstract Good Observable Class}

In \texttt{ModeStateObservableClass.lean}, Lean packages the good side on the
shared \texttt{ModeState}: nonnegative summable mode weights.  For that whole
class, Lean proves:
\begin{enumerate}[leftmargin=1.5em]
  \item continuity of the weighted radius observable;
  \item continuity of the derived statistic, cutoff potential, and Cole--Hopf
  datum;
  \item convergence of all of these along the existing truncations
  \(\texttt{truncateModes}\);
  \item recovery of the earlier geometric radius model as a special case.
\end{enumerate}

So the good side is no longer just one geometric example.  It is now an
abstract summable observable class on the shared mode space.

\subsection{Abstract Tail-Spike No-Go Theorem}

In \texttt{ModeStateTailSensitivityObstruction.lean}, Lean packages the bad
side abstractly too.  If an observable is constant on every remote
single-spike state \(\texttt{tailSpikeMode}(N)\) but takes a different value at
\(\texttt{modeZero}\), Lean proves:
\begin{enumerate}[leftmargin=1.5em]
  \item the raw observable is not continuous at \(\texttt{modeZero}\);
  \item the same obstruction lifts to cutoff potentials;
  \item the same obstruction lifts further to Cole--Hopf data.
\end{enumerate}

So the bad side is no longer tied only to the one hand-built
\texttt{tailPresenceRadius}.  It is now a reusable theorem pattern.

\subsection{First Concrete Chart Fork}

In \texttt{ModeStateChartFork.lean}, Lean forces the first explicit chart-style
fork on the shared mode space:
\begin{enumerate}[leftmargin=1.5em]
  \item the soft support observable
  \(t \mapsto |t|/(1+|t|)\), combined with summable mode weights, lies on the
  good side and yields a continuous truncation-friendly chart radius, statistic,
  cutoff potential, and Cole--Hopf datum;
  \item the hard support observable, which only asks whether any coefficient is
  nonzero, lies on the bad side and fails continuity at \(\texttt{modeZero}\),
  together with its cutoff and Cole--Hopf lifts.
\end{enumerate}

This is the first genuine theorem-level NS fork in Lean:
\begin{quote}
soft support belongs to the good summable side, hard support belongs to the
bad tail-sensitive side.
\end{quote}

\subsection{Sobolev-Style Chart Fork}

In \texttt{ModeStateSobolevChartFork.lean}, Lean pushes that fork one step
closer to manuscript shape by amplifying each mode by the polynomial factor
\(((n+1)^m)\) before applying the chart cutoff:
\begin{enumerate}[leftmargin=1.5em]
  \item the soft Sobolev observable, built from the bounded scalar
  \(t \mapsto |t|/(1+|t|)\) together with a concrete inverse-square envelope,
  is continuous and truncation-friendly;
  \item the corresponding hard Sobolev observable, which only asks whether any
  amplified mode is nonzero, is proved equal to the earlier hard support
  detector and is therefore not continuous at \(\texttt{modeZero}\), together
  with its cutoff and Cole--Hopf lifts.
\end{enumerate}

So the current Lean fork is now sharper:
\begin{quote}
even after Sobolev-style high-frequency amplification, the soft bounded version
lands on the good summable side, while the hard support version still lands on
the bad side.
\end{quote}

\subsection{Energy-Style Chart Fork}

In \texttt{ModeStateEnergyChartFork.lean}, Lean pushes the same shared-state
fork one step further toward an energy-style observable:
\begin{enumerate}[leftmargin=1.5em]
  \item the soft energy observable squares the Sobolev-amplified modes and
  then applies a compensating inverse-power envelope; Lean proves that this
  radius is continuous and truncation-friendly, and so are its cutoff and
  Cole--Hopf lifts;
  \item the hard energy observable asks whether any amplified squared mode
  crosses the threshold \(1\); Lean proves that this detector is stable on the
  remote spike family, disagrees at \(\texttt{modeZero}\), and is therefore not
  continuous, again together with its cutoff and Cole--Hopf lifts.
\end{enumerate}

So the current shared-mode-space fork is now sharper:
\begin{quote}
if the chart observable keeps enough inverse-power compensation to remain
summable after Sobolev amplification and squaring, Lean puts it on the good
approximation-friendly side; if it turns into a hard energy threshold on those
amplified modes, Lean puts it on the bad tail-sensitive side.
\end{quote}

\subsection{Log-Energy Chart Fork}

In \texttt{ModeStateLogEnergyChartFork.lean}, Lean tests a smoother radial
profile on the same shared mode space by replacing raw squared Sobolev energy by
\(\log(1+\mathrm{energy})\):
\begin{enumerate}[leftmargin=1.5em]
  \item the soft log-energy observable is continuous and truncation-friendly;
  the key formal domination is \(\log(1+u) \le u\), which lets Lean compare the
  log-energy terms to the already-controlled soft energy terms;
  \item the hard-side analogue asking whether \(\log(1+\mathrm{energy}) > 0\)
  still lands on the bad side: it is stable on the remote spike family,
  disagrees at \(\texttt{modeZero}\), and therefore its cutoff and Cole--Hopf
  lifts are also not continuous.
\end{enumerate}

So the shared-mode-space fork is now sharper again:
\begin{quote}
even replacing raw energy by the smoother profile \(\log(1+\mathrm{energy})\)
does not escape the same dichotomy: the summably compensated soft version stays
on the good side, while the hard detection version stays on the bad side.
\end{quote}

\subsection{Generic Energy-Profile Fork}

In \texttt{ModeStateEnergyProfileFork.lean}, Lean abstracts the preceding
examples into one theorem pattern.  An \texttt{EnergyProfile} is a globally
continuous radial profile \(\rho : \mathbb{R} \to \mathbb{R}\) satisfying, on
nonnegative inputs:
\begin{enumerate}[leftmargin=1.5em]
  \item \(0 \le \rho(u)\),
  \item \(\rho(u) \le u\),
  \item \(\rho(0)=0\),
  \item \(u>0 \Rightarrow \rho(u)>0\).
\end{enumerate}

For every such profile, Lean proves:
\begin{enumerate}[leftmargin=1.5em]
  \item the summably compensated soft radius is continuous and
  truncation-friendly;
  \item the induced hard detector agrees with the raw hard Sobolev support
  detector;
  \item that hard detector is tail-spike stable and therefore not
  continuous at \(\texttt{modeZero}\), together with its cutoff and Cole--Hopf
  lifts.
\end{enumerate}

The file also includes canonical instances:
\begin{enumerate}[leftmargin=1.5em]
  \item the identity profile, recovering the raw soft energy radius;
  \item the clipped log profile \(u \mapsto \log(1+\max(u,0))\), recovering the
  soft and hard log-energy constructions on the nonnegative energy inputs that
  actually occur;
  \item a bounded rational saturation profile
  \(u \mapsto \max(u,0)/(1+\max(u,0))\), which is closer to a chart-cutoff
  saturation shape and whose hard detector Lean proves collapses back to the
  hard Sobolev support detector;
  \item a smooth exponential saturation profile
  \(u \mapsto 1-\exp(-\max(u,0))\), again with the same hard-detector collapse;
  \item a smooth-transition profile
  \(u \mapsto u \cdot \mathrm{smoothTransition}(u)\), using mathlib's
  \texttt{Real.smoothTransition}, again with the same hard-detector collapse;
  \item more strongly, a positive-parameter family
  \(u \mapsto u \cdot \mathrm{smoothTransition}(cu)\) for \(c>0\), all with the
  same hard-detector collapse.
\end{enumerate}

So the current shared-mode-space NS fork is now theorem-level rather than
example-level:
\begin{quote}
every continuous radial profile that stays nonnegative, vanishes only at zero,
and is dominated by the identity on nonnegative inputs lands on the same
good/bad fork.
\end{quote}

\subsection{Finite Identity-Energy Algebra}

In \texttt{IdentityEnergy.lean}, Lean formalizes a finite
carr\'e-du-champ-style energy
\[
\gamma(D) = \sum_i D_i^2
\]
and proves:
\begin{enumerate}[leftmargin=1.5em]
  \item basic positivity and scaling laws for \(\gamma\);
  \item the log-gradient comparison used in the identity-only
  Cole--Hopf calculation;
  \item the finite Cauchy--Schwarz frame-evaluation bound;
  \item the resulting pointwise vorticity-style estimate after push-down
  through a curl frame.
\end{enumerate}

So the finite algebraic spine of the identity-only bound is already Lean-backed.

\subsection{Abstract Cole--Hopf/Vorticity Interface}

In \texttt{ColeHopfInterface.lean}, the finite identity-energy algebra is
packaged into an abstract structure carrying:
\begin{enumerate}[leftmargin=1.5em]
  \item uniform positivity of \(\Phi(t)\);
  \item a uniform identity-energy bound on its directional derivatives;
  \item a uniform frame-curl bound.
\end{enumerate}

From those hypotheses, Lean derives:
\begin{enumerate}[leftmargin=1.5em]
  \item a bound on the log-potential coefficients
  \(W(t) = -2\nu \log \Phi(t)\);
  \item a pointwise abstract vorticity bound;
  \item the uniform-in-time version of the same estimate.
\end{enumerate}

\subsection{Kernel-Semigroup Wrapper}

In \texttt{KernelSemigroupInterface.lean}, the same identity-only
Cole--Hopf/vorticity algebra is wrapped around a local Markov-kernel semigroup.
This does not add new analytic content.  It says that if a future NS evolution
is encoded as a local kernel semigroup, then the existing identity-energy
bridge can be reused unchanged.

\section{Current Boundary}

The following are \emph{not} yet formalized in Lean in the present NS lane:
\begin{enumerate}[leftmargin=1.5em]
  \item the manuscript's specific chart-cutoff datum on
  \(\mathrm{SDiff}^m\);
  \item the finite-mode truncation / approximation theorem actually needed for
  that datum;
  \item the local hypotheses \((H1)\), \((H2)\);
  \item the manuscript's claimed identity-only propagation lemma;
  \item the BKM-closing bridge for the full infinite-dimensional argument.
\end{enumerate}

So the honest current status is:
\begin{quote}
Lean has formalized the generic interface shell that a proof attempt would need,
but not the manuscript-specific analytic instantiations.
\end{quote}

\section{Verification}

The currently relevant NS modules were checked with:
\begin{verbatim}
lake build \
  Mettapedia.FluidDynamics.NavierStokes.ModeStateEnergyProfileFork \
  Mettapedia.FluidDynamics.NavierStokes.ModeStateLogEnergyChartFork \
  Mettapedia.FluidDynamics.NavierStokes.ModeStateEnergyChartFork \
  Mettapedia.FluidDynamics.NavierStokes.ModeStateSobolevChartFork \
  Mettapedia.FluidDynamics.NavierStokes.ModeStateChartFork \
  Mettapedia.FluidDynamics.NavierStokes.ModeStateTailSensitivityObstruction \
  Mettapedia.FluidDynamics.NavierStokes.ModeStateObservableClass \
  Mettapedia.FluidDynamics.NavierStokes.ModeStateProductTopologyObstruction \
  Mettapedia.FluidDynamics.NavierStokes.ModeStateProductTopologyContinuity \
  Mettapedia.FluidDynamics.NavierStokes.ProductTopologyL1Countermodel \
  Mettapedia.FluidDynamics.NavierStokes.GeometricModeApproximation \
  Mettapedia.FluidDynamics.NavierStokes.ChartApproximationModel \
  Mettapedia.FluidDynamics.NavierStokes.IdentityEnergy \
  Mettapedia.FluidDynamics.NavierStokes.WeakTopologyCountermodel \
  Mettapedia.FluidDynamics.NavierStokes.KernelSemigroupInterface \
  Mettapedia.FluidDynamics.NavierStokes.ApproximationInterface \
  Mettapedia.FluidDynamics.NavierStokes.CutoffContinuity \
  Mettapedia.FluidDynamics.NavierStokes.ColeHopfTopology \
  Mettapedia.FluidDynamics.NavierStokes.ColeHopfInterface
\end{verbatim}

\section{Next Lean Target}

The natural next formal step is no longer more prose.  It is to instantiate one
of the current abstract fork with manuscript-specific data:
\begin{enumerate}[leftmargin=1.5em]
  \item formalize a manuscript-shaped chart observable on the shared mode space
  whose shape is \emph{closer} than the current Sobolev-style soft/hard toy,
  and prove whether it belongs to the good summable class or the bad
  tail-sensitive class;
  \item if it lands on the good side, formalize the approximation/truncation
  theorem needed to feed \texttt{ApproximationInterface.lean};
  \item formalize the manuscript's local identity-energy bridge inside the
  abstract shell already provided by \texttt{IdentityEnergy.lean} and
  \texttt{ColeHopfInterface.lean}.
\end{enumerate}

\end{document}
